\section{Discussion}

\subsection{Interpretation of Results}

\paragraph{Cost versus aliasing} The central finding is that, in this implementation, the fidelity gained from anti-aliasing costs almost nothing in computation. PolyBLEP reduced the aliasing energy of a 1760\,Hz sawtooth by a large amount of 19.82\,dB (Figure~\ref{fig:alias}), and yet the band-limited oscillators ran within only about 1 to 2\,$\mu$s per block of the naive ones (Table~\ref{tab:cpu_osc}). The triangle with PolyBLAMP was the most expensive, at 3.4\% above its naive counterpart, but even this used only 2.21\% of the real-time block budget. The measured sawtooth result, in which the band-limited core was slightly faster than the naive one, is not evidence that PolyBLEP is cheaper. It is more plausibly measurement variation of the order of 1\,$\mu$s, and the safe conclusion is that the two have indistinguishable cost at the resolution of these measurements. This agrees with the expectation from the algorithms, since PolyBLEP and PolyBLAMP add a small correction only in the few samples around each discontinuity, so the extra work per block is small\cite{PolyBLEP}\cite{PolyBLAMP}.

\paragraph{Filter cost.} The ZDF state variable filter was usually more expensive than the Biquad IIR, by 2.51\,$\mu$s on average (about 4.4\%) across LFO rates of 10 to 50\,Hz, which matches the expectation that it would use more cycles. The difference is small, and it is smaller than the variation between repeated measurements would plausibly be. The IIR cost was almost constant across LFO rates, whereas the ZDF cost varied more. The 0\,Hz point is an exception, where ZDF was somehow 8.09\,$\mu$s \emph{cheaper} than IIR. 

We did not investigate this. One possible explanation is that the ZDF code does less work when the cutoff is not changing---perhaps values for the \texttt{std::tan} are cached, so that \texttt{UpdateCoefficients} does zero work on each sample, when the LFO is at 0Hz---but this is a hypothesis, not something the data show. A more mundane possibility is run-to-run variation, since each point is a single trace.

\paragraph{Stability.} The stability experiment produced the clearest qualitative contrast (Figure~\ref{fig:stability}). Under an LFO rate of 150\,Hz, a base cutoff of 57\,Hz and 100\% resonance, the Biquad IIR output saturated repeatedly at the limits of the output range, whereas the ZDF output remained bounded after an initial transient. This is consistent with the expectation that a direct-form biquad whose coefficients change at audio rate can inject energy, because its internal state was accumulated under coefficients that have since changed, while a ZDF structure that solves the feedback loop explicitly tolerates such changes better\cite{va_filter}. Because the output is clipped to $[-1,1]$, the recording shows that the IIR output diverged but cannot show how far the internal state had grown, and it provides no stability margin. The honest summary is therefore that the ZDF filter remained stable \emph{under the conditions tested}, but not that it is unconditionally stable; there could be a situation where both the ZDF and biquad filters are unstable.

\paragraph{The trade-off} Combining the three experiments answers the question posed in the title of this report; the measured costs of both anti-aliasing and ZDF filtering are small compared with the real-time budget, and the benefits are large. Specifically, a simple implementation of PolyBLEP and PolyBLAMP produced a reduction in aliasing of almost 20\,dB, and the ZDF method produced a filter that stays usable under modulation that makes the cheaper alternative unusable. In this setting there is little tension between fidelity and efficiency. The cheaper algorithms were not meaningfully faster, so choosing them saves almost nothing---in the context of a browser-based synthesizer---while costing obvious audible quality.

\subsection{Feasibility in the Browser}

The results also address the browser-specific part of the research question. Every configuration measured used no more than about 2.3\% of the 2.67\,ms block budget at 48\,kHz, so these algorithms are clearly practical in WebAssembly inside an AudioWorklet for a single voice. Two caveats apply. First, the measured times include fixed browser overhead around each call, not just the synthesizer, so the cost of an additional voice would be lower than the full per-block figure. Second, scaling to many voices has not been tested. A simple extrapolation from about 60\,$\mu$s per block would suggest a large number of voices, but this should not be relied on, since browser scheduling variation, memory behaviour and worst-case spikes (which the mean hides) are likely what causes dropouts in 
practice\footnote{\url{https://cprimozic.net/blog/webaudio-audioworklet-optimization/}}. The development experience also showed that the main difficulty in this environment was not DSP cost but architecture: the separation of the UI and audio threads, with parameters passed by message, constrained how modulation and automation could be implemented.

\subsection{Limitations}

This report answers parts of the research question more fully than others, and the following limitations should be kept in mind.
\begin{itemize}
  \item \textbf{Single measurements} Each CPU figure comes from one trace, with no repetition. Differences of a few microseconds, including the sawtooth result and the 0\,Hz anomaly, are within plausible variation. Reporting block times over repeated runs would make the cost comparisons much firmer.
  \item \textbf{Aliasing coverage} The quantitative alias measure was applied to the sawtooth at a single pitch (1760\,Hz). The PolyBLEP square and PolyBLAMP triangle were not assessed with the same measure. The metric is also an absolute energy and not a ratio to the harmonic energy, so it depends on both recordings having the same level.
  \item \textbf{Informal stability test} The stability experiment was a stress test at one extreme setting, not the automated sweep with automatic detection of non-finite values and clicks that was originally planned. It shows qualitative divergence in one filter but does not compare stability across filter modes (lowpass, highpass, bandpass), nor does it quantify margins.
  \item \textbf{Filter response accuracy} The accuracy of the filters' frequency responses was not measured, so that part of the audio fidelity question remains open.
  \item \textbf{Scope} The original aim of a multi-voice, multi-instrument synthesizer was not achieved. All experiments use a single voice, and some parameters (such as LFO rate, sustain and gain) are global rather than per voice.
\end{itemize}

\subsection{Relation to Previous Work}
\noindent
The results are consistent with the well-established account in the literature that PolyBLEP-style correction suppresses aliasing at low cost\cite{PolyBLEP}, and that zero-delay-feedback filters behave better than direct-form structures under fast parameter changes\cite{va_filter}. The contribution of this work is not a new algorithm but a controlled comparison of these methods within a single browser-based system, measured in the environment where they would be used, with a real-time budget that is much tighter and less predictable than in a native plugin.

\subsection{Future Work}
\noindent
The most valuable extensions would be: repeating the CPU measurements and reporting spreads and worst-case times; applying the alias measure to all waveforms and several pitches; replacing the informal stability test with the automated sweep, including detection of non-finite output and clicks, across all filter modes; measuring filter response accuracy against an analogue prototype; and testing polyphony to find the point at which real-time reliability is lost in the browser.